\section{Experiments}

In this section, we conduct a series of experiments to demonstrate the effectiveness of MAP across various tasks, including commonsense reasoning, natural language understanding, and subject-driven generation tasks.
In the following subsections, we provide detailed descriptions of each task and report the corresponding performance achieved by MAP.
All LoMAP runs use standard LoRA initialization, $\epsilon=10^{-6}$, $\alpha_0=\|\mathbf{W}\|_F$, and $\beta_0=1$.
The Frobenius normalization denominators are optimized jointly with the low-rank matrices rather than detached.


\input{sec/tables/nlu}



\subsection{Commonsense Reasoning}

\subsubsection{Task, Model, and Baselines}

The commonsense reasoning evaluation includes eight diverse benchmarks, each associated with a specific dataset: BoolQ \cite{clark2019boolq}, PIQA \cite{bisk2020piqa}, SIQA \cite{sap2019socialiqa}, HellaSwag \cite{zellers2019hellaswag}, WinoGrande \cite{sakaguchi2021winogrande}, ARC-e, ARC-c \cite{clark2018thinkarce}, and OpenBookQA (OBQA) \cite{mihaylov2018canobqa}. 
Following the protocol proposed in \cite{hu2023llm}, we consolidate the training splits of all benchmarks into a unified dataset, referred to as Commonsense170K, and evaluate model performance on the test sets of each benchmark individually.
We fine-tune LLaMA-7B \cite{touvron2023llama} and LLaMA3-8B \cite{llama3modelcard} on this target task. 

We compare LoMAP with several baselines including Prefix \cite{li2021prefix}, Series \cite{houlsby2019parameter}, Parallel \cite{he2021towards}, LoRA \cite{hu2021lora}, AdaLoRA \cite{zhang2022adaptive}, FLoRA \cite{si2024flora}, DoRA \cite{liu2024dora}, PiSSA \cite{meng2024pissa}, and MiLoRA \cite{wang2024milora}.
In addition, we include comparisons with ChatGPT (gpt-3.5-turbo) by leveraging its zero-shot Chain-of-Thought reasoning capabilities, as outlined in \cite{wei2022chain}.
All the experiments are conducted using NVIDIA A100 GPUs. 
The hyper-parameters are
shown in Table \ref{tab: cr detail}.


\begin{table}[!ht]

\centering
\renewcommand\arraystretch{1}
\setlength{\tabcolsep}{3mm}
\caption{Hyper-parameter settings of LoMAP on commonsense reasoning task.}
\resizebox{\linewidth}{!}{
\begin{tabular}{c | c c c c | c c} 

\toprule

\textbf{Settings} & \multicolumn{4}{c|}{LLaMA-7B} & \multicolumn{2}{c}{LLaMA3-8B}\\ 

\midrule


Rank $r$ & 4 & 8 & 16 & 32 & 16 & 32 \\ \midrule

$\alpha$ & 32 & 64 & 32 & 64 & 32 & 64 \\ \midrule


LR ($10^{-4}$) & 3 & 3 & 2 & 3 & 3 & 3 \\ \midrule

LR Scheduler & \multicolumn{6}{c}{Linear} \\ \midrule

Dropout & \multicolumn{6}{c}{0.05} \\ \midrule

Optimizer & \multicolumn{6}{c}{AdamW} \\ \midrule

Batch size & \multicolumn{6}{c}{16} \\ \midrule

Warmup Steps & \multicolumn{6}{c}{100} \\ \midrule

Epochs & \multicolumn{6}{c}{3} \\ \midrule

Where & \multicolumn{6}{c}{Q, K, V, Up, Down} \\

\bottomrule

\end{tabular}}
\label{tab: cr detail}
\end{table}





\begin{table*}
\centering
\renewcommand\arraystretch{0.9}
\setlength{\tabcolsep}{4mm}
\caption{Hyper-parameter settings of LoMAP on NLU task.}
\resizebox{\linewidth}{!}{
\begin{tabular}{c | c c c c c c c c} 

\toprule

Hyper-parameter & MNLI & SST-2 & CoLA & QQP & QNLI & RTE & MRPC & STS-B\\ 

\midrule

Optimizer & \multicolumn{8}{c}{AdamW} \\ \midrule

Warmup Ratio & \multicolumn{8}{c}{0.1} \\ \midrule

LR schedule & \multicolumn{8}{c}{Linear} \\  \midrule

Rank $r$ & \multicolumn{8}{c}{2 \& 8}\\ \midrule

LoRA alpha & \multicolumn{8}{c}{4 \& 16} \\ \midrule

Max Seq. Len. & 256 & 128 & 64 & 320 & 512 & 320 & 320 & 128 \\ \midrule

Batch Size & 32 & 32 & 32 & 32 & 32 & 32 & 32 & 32 \\ \midrule

Learning Rate & 5e-4 & 8e-4 & 8e-4 & 1e-3 & 5e-4 & 1.2e-3 & 1e-3 & 5e-4 \\ \midrule

Epochs & 12 & 24 & 25 & 5 & 5 & 50 & 30 & 25  \\ 

\bottomrule

\end{tabular}}
\label{tab: nlu detail}
\end{table*}




\subsubsection{Experimental Results}

Table \ref{tab:results of commonsense} presents the evaluation results of various PEFT methods on commonsense reasoning benchmarks under multiple model backbones and rank settings. 
We observe several notable trends:

First, LoMAP consistently outperforms all reported PEFT baselines under comparable parameter budgets.
For instance, under the LLaMA-7B backbone with $r=16$, LoMAP achieves an average accuracy of 77.9, substantially improving over LoRA (70.9); no DoRA result is available at this rank setting.
At $r=32$, LoMAP further reaches 78.8, surpassing DoRA (78.4) and LoRA-Dash (78.4) under similar parameter budgets.
Notably, even at lower ranks (e.g., $r=4$ or $r=8$), LoMAP yields strong performance, achieving 75.9 at $r=4$ and 76.4 at $r=8$, clearly outperforming the reported AdaLoRA, FLoRA, and DoRA baselines where available.
Second, LoMAP demonstrates excellent scalability across model sizes. 
When evaluated on the larger LLaMA3-8B model, LoMAP achieves the highest accuracy of 85.8 at $r=32$, surpassing strong baselines such as DoRA (85.2), and AdaLoRA (82.1). 
This indicates that LoMAP remains effective even when scaling to larger models and more complex reasoning tasks.
These results collectively validate the effectiveness and generality of LoMAP as a principled and scalable improvement over existing low-rank adaptation methods.



\subsection{Natural Language Understanding}


\subsubsection{Task, Model, and Baselines}

For the Natural Language Understanding (NLU) task, we use the General Language Understanding Evaluation (GLUE) benchmark \cite{wang2018glue}, which evaluates models across various tasks. 
The benchmark includes two sentence classification tasks: CoLA \cite{warstadt2019neural} and SST-2 \cite{socher2013recursive}, three tasks related to similarity and paraphrasing: MRPC \cite{dolan2005automatically}, QQP \cite{wang2018glue}, and STS-B \cite{cer2017semeval}, as well as three natural language inference tasks: MNLI \cite{williams2017broad}, QNLI \cite{rajpurkar2016squad}, and RTE \cite{dagan2005pascal,bar2006second,giampiccolo2007third,bentivogli2009fifth}. 
We fine-tune the DeBERTaV3-base and DeBERTaV3-large \cite{he2021debertav3} models on these tasks. 
The hyper-parameter settings are shown in Table \ref{tab: nlu detail}.


In addition to LoRA, DoRA and Series methods, we also include PAdapter in our comparisons.
Series introduces adapter modules at the interface between the self-attention and FFN blocks, incorporating them with residual connections to preserve model flow. In contrast, PAdapter employs a more streamlined design by attaching adapters solely after the FFN and LayerNorm layers.





\subsubsection{Experimental Results}


Table \ref{tab: deberta results} presents the GLUE benchmark results with DeBERTaV3-base and DeBERTaV3-large under various PEFT settings. 
Across both model sizes and rank configurations, LoMAP consistently delivers superior performance.

Under the DeBERTaV3-base setting, LoMAP exhibits clear advantages even in low-rank regimes.
At rank $r=2$, it achieves an average score of 89.49, outperforming LoRA and Padapter, while maintaining comparable parameter efficiency. 
As the rank increases to $r{=}8$, LoMAP continues to lead, reaching 89.64, surpassing LoRA’s 88.27 by a substantial margin.
For the larger DeBERTaV3-large model, the benefits of LoMAP remain prominent. 
With $r{=}8$, it achieves 91.18 average score, closing the gap with full fine-tuning (91.36) while requiring less than 1\% of the trainable parameters. 
This demonstrates LoMAP’s strong capacity to scale to more expressive architectures without sacrificing efficiency.
These results confirm the practicality of LoMAP for general-purpose language understanding and its potential to replace conventional fine-tuning in resource-constrained settings.


\subsection{Subject-driven Generation}


\begin{figure}[t]
    \centering
    \includegraphics[width=\linewidth]{fig/sdg.pdf}
    \caption{Comparison of generated images from LoRA and LoMAP on the subject-driven generation task. It is evident that LoMAP consistently produces images that better reflect both the input subjects and the intended prompts compared to standard LoRA.}
    \label{fig: sdg}
\end{figure}





\subsubsection{Task, Model, and Baselines}

In this experiment, we fine-tune text-to-image diffusion models for subject-driven image generation, following the setup proposed in DreamBooth \cite{ruiz2023dreambooth}.
The goal is to synthesize images that faithfully reflect a specific subject, given only a few reference examples. 
To achieve this, we fine-tune a text-to-image model using image-text pairs in which the subject is denoted by a unique identifier (e.g., ``A photo of a [V] cat'').
After fine-tuning, this identifier is embedded into new prompts to guide image generation specific to the learned subject.

We use the SDXL model \cite{podell2023sdxl} as the backbone and apply both LoRA and LoMAP as adaptation techniques. 
The model is trained with a learning rate of 1e-4, a batch size of 4, and for 500 steps on a single 80GB A100 GPU, which takes approximately 26 minutes. 
Image generation is conducted using 50 inference steps per prompt, with each synthesis taking around 10 seconds.
All experiments are conducted using the official DreamBooth dataset \cite{ruiz2023dreambooth}.








\subsubsection{Experimental Results}


As illustrated in Fig. \ref{fig: sdg}, the qualitative results highlight that LoMAP yields images with greater subject fidelity compared to standard LoRA. 
In particular, LoRA’s generated samples—for example, those depicting a dog or a teapot—often diverge noticeably from the reference images.
In contrast, LoMAP consistently preserves key subject attributes, producing visuals more closely aligned with the original exemplars.
In addition, LoMAP demonstrates strong semantic alignment with complex prompts, accurately interpreting and visually rendering fine-grained concepts such as \textit{cobblestone} or \textit{white rug}.
This highlights LoMAP’s capacity to effectively disentangle and integrate subject- and prompt-specific information during synthesis.
